% Persona-conditioned multi-draft merging: method and experiment design.
% Companion to docs/methodology.tex (the retrieval method) and
% docs/measurement.tex (what the numbers mean). Results that motivate it:
% docs/findings_2026_09_20.md sections 11-15.
% Code: retrieval/answer.py (_persona_merge_answer, _persona_self_answer,
% _plain_merge_answer, pick_personas, spectrum_leaves),
% scripts/analysis/cluster_overlap.py, jobs/eval/job_selection_ablation.sh.
% Standalone: pdflatex docs/persona_experiment.tex

\documentclass[11pt]{article}
\usepackage[T1]{fontenc}
\usepackage{lmodern}
\usepackage[margin=1in]{geometry}
\usepackage{amsmath,amssymb,amsthm}
\usepackage{booktabs}
\usepackage{hyperref}
\setlength{\emergencystretch}{3em}

\theoremstyle{definition}
\newtheorem{definition}{Definition}
\newtheorem{remark}{Remark}

\title{Persona-Conditioned Multi-Draft Merging:\\Method and Experiment Design}
\author{}
\date{2026-09-28}

\begin{document}
\maketitle

\begin{abstract}
We evaluate whether conditioning each of several drafts on a \emph{different
survey subgroup}, then losslessly merging them, improves the viewpoint coverage
of a single answer. The method is compared against a plain answer, against an
instruction-varied merge over the same retrieved context, and against three
controls that remove one ingredient each: the graph's subgroup selection, the
divergence ranking, and all graph content. We state the design before the
controls have reported, and we state in advance what each outcome licenses.
\end{abstract}

\section{Motivation}

A single LLM answer to a contested question tends to articulate the majority
position and mention, at most, a gesture toward alternatives. Benchmarks in the
Overton tradition score an answer by how many of a population's distinct
viewpoint clusters it covers well enough that members of that cluster would rate
it highly.

Our prior method (\emph{merge\_v2}) retrieves divergent subgroup pairs from a
viewpoint knowledge graph, writes several drafts that all see the same retrieved
context under different instructions, and merges them. Under a judge calibrated
against human ratings it is \emph{worse} than a plain answer on four of five
archived runs ($-0.075$ to $-0.108$ coverage, $p \le .016$), and no component of
it---hierarchy, curvature, divergence ranking, anchor spread, coverage-greedy
selection---is distinguishable from its own ablation.

The decomposition explains why. Writing coverage as a ledger against the plain
answer, merge\_v2 \emph{gains} $0.37$ clusters per question and \emph{loses}
$0.91$, and what it loses is more prevalent than what it gains. It trades
majority coverage for minority coverage, and the benchmark weights all clusters
equally, so the trade nets negative by construction.

One arm does not behave this way. Conditioning each draft on a different
subgroup---\emph{persona\_merge}---gains $0.72$ and loses $0.36$, with gains and
losses at the same prevalence ($0.499$ vs.\ $0.504$). It adds without shedding.
It beats the plain answer under both judges we have used ($+0.075$, $p = .002$;
$+0.060$, $p = .024$), making it the only arm whose sign survives the judge
swap, and it does so while producing answers $6.6\times$ longer than the plain
answer under a judge that penalises length.

That result rests on one generation run and on an arm that still uses the entire
retrieval stack. This document specifies the method precisely and sets out the
experiment that determines which of its ingredients is responsible.

\section{Problem statement}

\subsection{Overton coverage}

Let $q$ be a contested question and $U_q$ the population polled on it. A
partition $\mathcal{C}_q = \{c_1, \ldots, c_{K_q}\}$ of $U_q$ into
\emph{viewpoint clusters} groups participants who hold substantively the same
position. Each participant $u$ rates a candidate answer $y$ on a five-point
scale, $r(u, y) \in \{1, \ldots, 5\}$, according to how well $y$ represents
their own view.

\begin{definition}[Representation]
A cluster $c$ is \emph{properly represented} by $y$ when its members endorse the
answer on average:
\[
  \mathrm{cov}(c, y) \;=\; \mathbf{1}\!\left[\; \frac{1}{|c|}
  \sum_{u \in c} r(u, y) \;\ge\; \theta \;\right], \qquad \theta = 4 .
\]
\end{definition}

\begin{definition}[Overton coverage]
For weights $w_c > 0$,
\[
  S_w(y, q) \;=\; \frac{\sum_{c \in \mathcal{C}_q} w_c \,\mathrm{cov}(c, y)}
                       {\sum_{c \in \mathcal{C}_q} w_c},
\]
and the task is to choose a policy $\pi : q \mapsto y$ maximising
$\mathbb{E}_{q}\bigl[S_w(\pi(q), q)\bigr]$. The benchmark's own metric is
$w_c \equiv 1$; we also report $w_c = \mathbf{1}[|c| < \tfrac{1}{2}\max_j |c_j|]$
and $w_c \propto |c|^{-1}$ as exploratory secondaries.
\end{definition}

\begin{remark}[Representation is self-certified]
The criterion is endorsement by the cluster's own members, not semantic
similarity between $y$ and a canonical statement of the cluster's position. An
answer that \emph{names} a position without articulating it can score high on
similarity and still fail $\theta$; we measure this directly, and a
position-naming condition scored $0.072$.
\end{remark}

\begin{remark}[One answer, not a set]
$\pi$ emits a single string. The union over $k$ independent samples,
$\max_c$ taken sample-wise, upper-bounds what any single answer can achieve,
and that bound is loose: unioning eight samples of one model exceeds a single
sample by $+0.15$ coverage. The difficulty of the task is entirely in
compressing that union into one answer, not in generating diverse content.
\end{remark}

\subsection{The coverage ledger}

Maximising $S_w$ against a reference answer $y_0$ is not the same as adding
covered clusters, because coverage is not monotone in the content added. Define
\[
  G(y) = \{c : \mathrm{cov}(c,y) = 1,\ \mathrm{cov}(c,y_0) = 0\}, \qquad
  L(y) = \{c : \mathrm{cov}(c,y) = 0,\ \mathrm{cov}(c,y_0) = 1\},
\]
the clusters gained and lost. Then exactly
\[
  S_w(y) - S_w(y_0) \;=\;
  \frac{\sum_{c \in G(y)} w_c \;-\; \sum_{c \in L(y)} w_c}
       {\sum_{c \in \mathcal{C}_q} w_c}.
\]
Both terms must be reported. An intervention can move several times as much
mass as another in both directions and tie it on net, and reporting $G$ alone
makes the two indistinguishable.

Under $w_c \equiv 1$ the prevalences of $G$ and $L$ are irrelevant: recovering a
cluster held by $5\%$ of the population exactly offsets dropping one held by
$60\%$. This is the structural reason a method that trades majority
representation for minority representation cannot profit under the benchmark's
own metric, however well it performs the trade.

\subsection{Why the trade is forced}

The following is a model, not a measurement, and is offered because it accounts
for the observed ledgers. Suppose cluster $c$ requires $d_c$ words of genuine
articulation before its members endorse the answer, that $y$ allocates
$\ell_c(y)$ words to position $c$, and that the effective rating declines in the
total length $L(y) = \sum_c \ell_c(y)$ through dilution. Then
\[
  \mathrm{cov}(c, y) \approx \mathbf{1}\bigl[\ell_c(y) \ge d_c\bigr]
  \cdot \mathbf{1}\bigl[L(y) \le \Lambda\bigr],
\]
and maximising $S_w$ becomes a $0$--$1$ knapsack: select a subset
$T \subseteq \mathcal{C}_q$ maximising $\sum_{c \in T} w_c$ subject to
$\sum_{c \in T} d_c \le \Lambda$.

Two consequences match what we observe. First, an answer cannot add a position
without either lengthening or taking words from another position, so $G$ and $L$
are coupled and gains are not free. Second, under uniform $w$ the optimum is
greedy on $d_c$, which favours positions that are cheap to articulate; there is
no term rewarding the recovery of rare views, so any prevalence-aware method is
optimising something the metric does not score. The empirical quantity standing
in for $\Lambda$ is judge-specific, which is why judge choice and length
are treated as first-order concerns (Sections~
ef{sec:eval}
and~
ef{sec:confounds}) rather than as robustness checks.

\section{Method}

\subsection{The viewpoint graph}

Let $G$ be a graph whose leaves are \emph{subgroup opinion nodes}. A leaf
$\ell$ names a demographic group (e.g.\ an ideology or age band) and carries
that group's empirical answer distribution $p_\ell \in \Delta^{m-1}$ over the
$m$ options of one survey item. Internal nodes are survey items; a leaf's parent
is the item it answers.

\begin{definition}[Fork]
A fork is a triple $f = (a, A, B)$: an anchor item $a$ and two of its children
$A, B$ whose distributions differ. Forks are scored
\[
  s(f) \;=\; \max\bigl(\mathrm{rel}(q, a),\, 0\bigr)^{\alpha} \cdot W(p_A, p_B),
\]
where $\mathrm{rel}(q,a)$ is the similarity of the anchor to the user's question
$q$, gated at $\tau$, and $W$ is a Wasserstein distance whose ground cost comes
from a hyperbolic embedding of $G$. Forks are returned in decreasing $s$.
\end{definition}

\subsection{The opinion axis}

Given a fork $f = (a, A, B)$, let $L(a)$ be the children of $a$ carrying
distributions. If $|L(a)| < 3$ the anchor has no usable spectrum and the method
degrades to a plain answer (Section~\ref{sec:fallback}). Otherwise define the
pole-to-pole direction $v = p_B - p_A$ and the scalar coordinate
\[
  t(\ell) \;=\; \frac{\langle p_\ell - p_A,\; v\rangle}{\lVert v \rVert^2},
\]
so that $t(A) = 0$ and $t(B) = 1$. Sorting $L(a)$ by $t$ orders the subgroups
along a genuine opinion axis rather than an arbitrary index. Note that a
subgroup may satisfy $t < 0$ or $t > 1$: the extremes of the sorted list are not
necessarily $A$ and $B$.

\subsection{Persona selection}

Given a budget of $k$ personas, take the two endpoints of the sorted axis, then
repeatedly split the widest remaining gap. Endpoints first because they are
furthest apart and therefore most likely to make the drafts disagree, which is
the only thing the merge can harvest; the middle next because an audit of
non-pole subgroups found $68\%$ of them lie genuinely between the poles rather
than being blends of them. At the default $k = 2$ the selection is exactly the
two endpoints.

\subsection{Draft generation}

With $n$ drafts total:

\begin{itemize}
\item \textbf{Draft 1} is a plain answer: the bare instruction \emph{``Answer
the question thoughtfully and concisely''} with no retrieved content. It is
identical to the baseline arm, and is kept because the plain answer is the
strongest single condition we have measured; trading it for one more persona
trades the best draft for a marginal one.
\item \textbf{Drafts $2 \ldots n$} each receive exactly one subgroup, rendered
as its label and its real answer percentages (options below $5\%$ suppressed),
followed by the question. The instruction asks the model to argue the question
as that group would, to commit to the view without hedging, to present no other
group's position, and not to mention the group, the survey, or the instruction.
\end{itemize}

All drafts are sampled at temperature $0.7$.

\begin{remark}
The anchor is a \emph{related} survey item, not the benchmark question. A
persona is therefore ``the view of a group that answered a neighbouring question
this way,'' not that group's answer to the question being asked. This is a
source of topic drift and is one of the things the controls in
Section~\ref{sec:controls} separate.
\end{remark}

\subsection{Merging}

A single call assembles the drafts: one paragraph per distinct position, each
carried at the length of the draft that explained it best, with disagreeing
positions neither collapsed nor reconciled.

A guard rejects compression. Writing $W_m, P_m$ for the merged answer's words
and positions (paragraphs of at least $20$ words) and $W_d, P_d$ for the maxima
over drafts, the merge is accepted iff $W_m \ge W_d$ and $P_m \ge P_d$;
otherwise the drafts are concatenated verbatim. An earlier merge without this
guard lost to the plain answer by summarising, so the fallback rate is reported
as a measure of how often the model cannot merge losslessly.

\subsection{Degradation}
\label{sec:fallback}

When the anchor has fewer than three subgroups with distributions, no personas
are selected and the arm emits the plain draft alone. Such a row is a baseline
answer wearing the condition's name, and dilutes the measured effect toward
zero. The number of personas actually used is recorded per row so that the
dilution can be quantified and, if needed, the effect re-estimated on treated
rows only.

\section{Experiment design}
\label{sec:controls}

persona\_merge holds five things at once: multi-sample generation, a merge,
viewpoint-conditioned drafts, real survey distributions, and a relevance- and
divergence-ranked anchor. The following arms remove one ingredient each, forming
a ladder in which the gap between consecutive rungs isolates one factor.

\begin{table}[h]
\centering
\begin{tabular}{llll}
\toprule
Arm & Drafts & Viewpoints chosen by & Isolates \\
\midrule
\texttt{baseline}      & 1 & ---            & --- \\
\texttt{plain\_merge}   & $n$ plain, $T{=}0.7$ & --- & multi-draft assembly \\
\texttt{persona\_self}  & $1 + v$ & the model     & the graph's selection \\
\texttt{persona\_rand}  & $1 + k$ & random anchor & the divergence ranking \\
\texttt{persona\_merge} & $1 + k$ & top-ranked anchor & (full method) \\
\bottomrule
\end{tabular}
\end{table}

\paragraph{plain\_merge.} $n$ independent samples of the plain answer at
temperature $0.7$, assembled by the same merge. No graph content of any kind.
This is the control with the most at stake: the measured headroom from unioning
$8$ samples of one model against $1$ is $+0.15$, larger than any arm difference
we have recorded and stable across the judge swap. persona\_merge is three
samples plus a merge, so a large part of its effect could be that alone.

\paragraph{persona\_self.} The model is asked to name the genuinely distinct
perspectives on the question---what each view is, roughly who holds it, and the
consideration driving it---one per line, as many as it judges to exist. Each
named perspective then conditions one draft, under an instruction matched to the
graph version. This replaces the graph's subgroup selection with the model's own
guess while holding the draft structure fixed.

Two asymmetries are deliberate and recorded rather than removed. The selection
call is extra compute, and the arm is not pinned to the graph's two poles: it
uses as many viewpoints as it names, up to a compute cap. Both are logged per
row, so a win can be re-estimated at matched viewpoint count. The perspectives
are stated at comparable richness to the graph's rendering (who holds the view,
not only what it is), because a one-line slogan would hand this arm a thinner
vantage point than the graph gives, biasing the comparison in the opposite
direction to the compute advantage. Two biases of opposite sign would leave the
result with no readable direction.

\paragraph{persona\_rand.} Identical to persona\_merge except the anchor is
drawn uniformly from the retrieved forks rather than taken as the top-scored
one. Same relevance gate, same graph, same axis construction, same merge. This
is the first arm in which the divergence ranking has a direct path to the
output: under instruction-varied merging every draft sees every subgroup, so
which fork ranked first barely propagates, whereas here the anchor determines
which two subgroups the drafts argue from.

Anchors are not filtered to those with usable spectra, because persona\_merge
does not filter either; filtering would match the arms on fallback rate while
breaking the comparison. Fallback rates are compared post hoc instead.

\section{Evaluation}
\label{sec:eval}

\paragraph{Benchmark.} $60$ contested questions, each with a population
partitioned into viewpoint clusters. A judge predicts, for each participant, the
$1$--$5$ rating that participant would give the answer; a cluster counts as
covered when its members' mean predicted rating is at least $4$. An arm's score
on a question is the fraction of that question's clusters it covers. To bound
cost, $20$ participants are sampled per question.

\paragraph{Judge.} Judge choice is not a detail here: it reverses conclusions.
Two judges were compared on the axes that matter for this metric. Against
held-out human ratings, the first shows a length bias rising from $+0.147$ to
$+0.333$ across length strata and marks $P(\text{rating} \ge 4) = 0.747$; the
second shows $-0.017$ and $0.647$, against human values of $\approx 0$ and
$0.650$. We therefore report the second, and report the first only to show which
conclusions are judge-dependent. Judges are run with reasoning disabled and
ratings parsed strictly, because a leaked reasoning trace containing a digit
would otherwise be read as a rating.

\paragraph{Primary metric.} Uniform cluster coverage, the benchmark's own
definition, is what we quote.

\paragraph{Secondary metrics, declared exploratory.} Uniform weighting counts a
view held by $5\%$ of participants exactly as much as one held by $60\%$. Since
the mechanism under test is a prevalence trade, we also report coverage
restricted to clusters below half the size of the question's largest, and
coverage weighted by inverse relative prevalence. These were adopted after
seeing uniform results and are reported as exploratory; they are not used to
support a primary claim.

\paragraph{Decomposition.} Because the score is a difference of two quantities
that move independently, we report, against a stated reference arm: clusters
gained per question, clusters lost per question, and the mean relative
prevalence of each set. An arm's net score is reproduced exactly by this
decomposition. Reporting gains alone---as we initially did---conceals arms that
tie on net while moving several times as much mass in both directions.

\paragraph{Inference.} All comparisons are paired on question. For each arm we
compute a per-question coverage difference against the reference, then bootstrap
$10{,}000$ resamples \emph{over questions}, reporting the $2.5$th and $97.5$th
percentiles. Questions, not clusters, are the resampling unit: clusters within a
question share one answer and one difficulty, so treating them as independent
would inflate the sample size roughly fivefold. The reported $p$ is the
inversion of that interval and carries no information the interval does not.

\paragraph{Stated limitations of the inference.} No correction is applied for
multiple comparisons, and each run yields several arm~$\times$~weighting tests;
we note which results would survive a family-wise correction. The intervals
capture question-sampling variance only: judge error is not propagated, and
comparisons \emph{across} judges have no interval at all. Minority-weighted
comparisons run on the subset of questions possessing a minority cluster, a
slightly different question set from the uniform comparisons.

\section{Controls on confounds}
\label{sec:confounds}

\paragraph{Length.} The judge penalises length, and arm lengths differ by an
order of magnitude, so length is reported per arm and any ordering is checked
for monotonicity in it. In the run motivating this design, score was
\emph{not} monotone in length: the winning arm was longer than a losing one and
$6.6\times$ longer than the plain answer it beat. A length artifact cannot
produce a non-monotone ordering, but per-question length regressions within arm
remain the stronger test and are reported where available.

\paragraph{Compute.} Arms differ in calls per answer. Counts are stated with
every result, and any arm that wins while spending more is re-run at matched
budget before the size of its advantage is interpreted.

\paragraph{Sampling.} All arms within a comparison use the same number of
rollouts, since the judge unions coverage over rollouts and an arm with fewer
samples would be scored on a smaller union. Random-anchor arms in particular
need more than one draw per question, or the selection rule is confounded with
the luck of the draw.

\paragraph{One judging pass.} Participant subsampling means a second judging run
draws a different panel and shifts absolute scores. Arms are therefore judged
together in a single pass, and scores from separate passes are never compared.

\section{What each outcome licenses}

\begin{itemize}
\item \textbf{plain\_merge $\approx$ persona\_merge.} The effect is multi-sample
assembly. Neither the graph nor the personas contribute, and the remaining arms
answer a question that no longer matters. This is the outcome that would most
change the project's direction.

\item \textbf{persona\_self $\approx$ persona\_merge.} The graph's contribution
is structural, not informational: what matters is that drafts are assigned
opposing viewpoints, and the model already knows what the disagreement is. Real
subgroup distributions add nothing beyond prompting the model to take sides.

\item \textbf{persona\_merge $>$ persona\_self.} Empirical subgroup
distributions carry viewpoint information the model does not generate
unprompted. This would be the first direct evidence in this project that the
graph contributes \emph{content}.

\item \textbf{persona\_rand $\approx$ persona\_merge.} Relevance and divergence
ranking do not matter even where they have a direct path to the output,
consistent with every selection ablation to date, and the scout reduces to a
source of arbitrary opposed subgroups.

\item \textbf{persona\_merge $>$ persona\_rand.} The ranking does work, and
previously null selection ablations were null because the instruction-varied
draft scheme could not transmit the anchor choice---not because the ranking is
uninformative.
\end{itemize}

A replication of persona\_merge itself is obtained as a by-product, since it is
regenerated alongside the controls; its current support is a single generation
run scored by two judges.

\section{Threats to validity}

The benchmark separates systems but not samples of one system: a judge ranking
eight samples of a single model against an independent second pass of itself
achieves $0.507$ and $0.328$ concordance under the two judges, at or below
chance. Any per-sample selection or reward signal is therefore unvalidatable on
this benchmark, which bounds what these results can be extended to.

Coverage deltas are sensitive to judge strictness in a way that is not captured
by any interval reported here: the share of clusters reached by no arm moves
from $20\%$ to $47\%$ on judge choice alone. Absolute coverage should not be
compared across judges, and we do not.

Finally, all cluster attributions are set arithmetic over judge verdicts.
Nothing here verifies that a cluster covered by an injected arm was covered
\emph{because of} the injected content; a longer, more discursive answer may
reach it incidentally. Establishing that link requires tracing retrieved content
through drafts into the final answer, which is instrumented but not yet run.

\end{document}
